\chapter{Literature Review}
\label{ch:literature_review}

% ─────────────────────────────────────────────────────────────────────────────
\section{Theoretical Foundations of Adaptive Learning}
\label{sec:lr_theory}
% ─────────────────────────────────────────────────────────────────────────────

\subsection{The 2-Sigma Problem and ITS}
\label{subsec:lr_2sigma}

\textcite{bloom1984} demonstrated that one-to-one human tutoring produces
learning outcomes two standard deviations above conventional instruction
(the ``2-sigma problem''). \textcite{vanlehn2011} partially resolves this
at scale: step-based ITS achieves $d = 0.76$, statistically indistinguishable
from human tutoring ($d = 0.79$). The critical determinant is interaction
granularity, not algorithmic sophistication---a finding that directly informs
the scaffolding architecture of this system. \textcite{chi2008} provide a
complementary account: the effectiveness of human tutors derives less from
domain expertise and more from eliciting and responding to student reasoning,
suggesting that ITS design should prioritise responsive feedback over
predictive accuracy.

\subsection{Zone of Proximal Development}
\label{subsec:lr_zpd}

Vygotsky's \parencite*{vygotsky1978} Zone of Proximal Development (ZPD)
defines the cognitive space between independent and scaffolded achievement.
Tasks below the ZPD produce boredom; tasks above it produce cognitive overload;
only tasks within it produce optimal learning. This system operationalises
the ZPD through IRT target probability $P(\theta) \approx 0.70$: the selected
item presents difficulty at which the learner has approximately 70\% probability
of success. The 13-KC prerequisite graph encodes the multi-dimensional
knowledge structure required to navigate Geography's integration of declarative,
procedural, and conceptual knowledge \parencite{park2004}.

\subsection{Mastery Learning and Bloom's Taxonomy}
\label{subsec:lr_mastery}

\textcite{bloom1984} proposed mastery learning: progression is gated by
demonstrated competence rather than time-on-task. \textcite{ritter2016} identify
80\% accuracy as the optimal threshold from large-scale Cognitive Tutor data;
this system implements $p(\text{Learned}) \geq 0.80$ per KC as the BKT mastery
gate. Bloom's revised taxonomy \parencite{anderson2001} provides three cognitive
levels (Remembering, Understanding, Applying); the adaptive engine enforces
escalation gates at $p(\text{Learned}) \geq 0.40$ (L1$\to$L2) and
$\geq 0.60$ (L2$\to$L3), ensuring declarative knowledge is consolidated before
analytical demands arise.

\subsection{Spaced Repetition and Cognitive Load}
\label{subsec:lr_spacing}

\textcite{cepeda2006} establish the spacing effect: distributed practice
consistently outperforms massed practice at delayed retention. This system
implements Successive Relearning \parencite{kornell2009} through BKT decay
detection: KCs whose $p(\text{Learned})$ has declined receive a $1.5\times$
information boost at next selection, prioritising re-presentation of
consolidating material. \textcite{sweller1988}'s cognitive load theory
informs the graduated hint system: each hint level reduces intrinsic load
while also reducing the evidential weight of a subsequent correct response
\parencite{koedinger2007}, a relationship formalised by the credit-factor
BKT modification.

% ─────────────────────────────────────────────────────────────────────────────
\section{Adaptive Algorithms: A Critical Comparative Analysis}
\label{sec:lr_algorithms}
% ─────────────────────────────────────────────────────────────────────────────

\textcite{matos2025} identify three dominant ITS algorithm families: Elo-based
rating systems, probabilistic models (IRT and BKT), and deep learning
approaches. Each is evaluated below against four project requirements:
per-KC mastery tracking, prerequisite enforcement, item discrimination
modelling, and interpretability.

\subsection{Elo Rating System}
\label{subsec:lr_elo}

Originally developed for chess player ranking \parencite{elo1978}, the Elo
system updates learner and item ratings after each response:

\begin{equation}
  R_i^{\text{new}} = R_i + K \cdot \left(\text{outcome} - P(\text{correct})\right)
  \label{eq:elo_update}
\end{equation}

\textcite{pelanek2016} reports approximately 92\% prediction accuracy after
25--30 interactions. However, Elo operates as a single global ability estimate:
it cannot model per-KC mastery states, cannot enforce prerequisite dependencies,
and provides no knowledge decay detection. These three limitations make Elo
unsuitable for a 13-KC curriculum with prerequisite structure, regardless of
its calibration simplicity.

\subsection{Item Response Theory (3PL)}
\label{subsec:lr_irt}

IRT \parencite{embretson2000} models response probability as a function of
latent ability. The three-parameter logistic (3PL) model \parencite{birnbaum1968}
is:

\begin{equation}
  P(X=1 \mid \theta) = c + \frac{1-c}{1 + e^{-a(\theta - b)}}
  \label{eq:irt_3pl}
\end{equation}

where $a$ is item discrimination, $b$ is difficulty, and $c$ is the
pseudo-guessing parameter (fixed at 0.25 for four-option items). The guessing
parameter is critical: Geography MCQ items carry a non-trivial guessing
probability, and ignoring it systematically overestimates learner ability
following correct responses. The principal limitation is calibration data
requirements: traditional 3PL estimation requires 500+ observations per item
\parencite{lord1980}, infeasible at student-project scale. This project
mitigates by fixing $c = 0.25$ as an equality constraint and pre-specifying
a Rasch (1PL) fallback if RMSEA $> 0.10$ \parencite{chalmers2012}. IRT has no
native mechanism for per-KC longitudinal tracking, requiring combination with
a knowledge state model.

\subsection{Bayesian Knowledge Tracing}
\label{subsec:lr_bkt}

BKT \parencite{corbett1995} conceptualises learning as a hidden Markov process
with binary latent states: a learner either knows or does not know a KC.
Four parameters govern each KC: $p(L_0)$ (prior knowledge), $p(T)$ (learning
transition), $p(S)$ (slip), and $p(G)$ (guess). After each observation the
posterior updates via Bayes' theorem:

\begin{equation}
  p(L_n \mid \text{obs}) = \frac{P(\text{obs} \mid L_n) \cdot p(L_n)}{P(\text{obs} \mid L_n) \cdot p(L_n) + P(\text{obs} \mid \neg L_n) \cdot (1 - p(L_n))}
  \label{eq:bkt_update}
\end{equation}

with learning transition:

\begin{equation}
  p(L_{n+1}) = p(L_n \mid \text{obs}) + \bigl(1 - p(L_n \mid \text{obs})\bigr) \cdot p(T)
  \label{eq:bkt_transition}
\end{equation}

BKT's critical strength for this project is per-KC resolution, enabling
mastery-gated progression. \textcite{beck2007} identify a parameter
identifiability problem for sparse data; this project mitigates it with
informed priors and domain-standard $p(S)$/$p(G)$ values following
\textcite{pardos2011}.

\textbf{Credit factor modification.} Standard BKT treats all correct responses
as equally informative. \textcite{koedinger2007} demonstrate that correct
responses following worked examples are substantially less informative than
unaided correct responses. This project modifies the learning transition by
credit factor $\text{cf} \in \{1.0, 0.9, 0.7, 0.5\}$ depending on scaffolding
level accessed:

\begin{equation}
  p(L_{n+1})^* = p(L_n \mid \text{obs}) + \bigl(1 - p(L_n \mid \text{obs})\bigr) \cdot p(T) \cdot \text{cf}
  \label{eq:bkt_credit}
\end{equation}

\subsection{Deep Knowledge Tracing}
\label{subsec:lr_dkt}

\textcite{piech2015} introduced DKT, replacing BKT's probabilistic model with
an LSTM network achieving AUC improvements of approximately 6\% over BKT.
\textcite{khajah2016} demonstrate that this advantage largely disappears when
BKT is extended with temporal features. More critically, DKT sacrifices
interpretability entirely: hidden state vectors carry no psychologically
meaningful correspondence, and adaptive decisions cannot be traced to specific
knowledge estimates. \textcite{matos2025} confirm that in small-scale
educational deployments, interpretable probabilistic systems remain more
appropriate. DKT is unsuitable for a project requiring full accountability
of every adaptive decision.

\subsection{Hybrid IRT + BKT: Rationale and Prior Work}
\label{subsec:lr_hybrid}

The hybrid approach integrates IRT's item-level psychometric precision with
BKT's per-KC knowledge tracking. Each algorithm addresses the other's primary
limitation: BKT provides no item discrimination; IRT provides no per-KC mastery
tracking. \textcite{kaser2017} confirm across 1.2~million student--problem
interactions that the value of hybrid architectures lies not in marginal
accuracy improvement but in simultaneous modelling of item properties and
temporal knowledge dynamics. \textcite{desmarais2012} review 40 learner
modelling systems and conclude that hybrid probabilistic approaches consistently
outperform single-algorithm systems in structured multi-KC curricula.
Table~\ref{tab:algo_comparison} summarises the algorithm selection rationale.

\begin{table}[H]
\centering
\caption{Comparative evaluation of adaptive algorithms against project requirements.}
\label{tab:algo_comparison}
\renewcommand{\arraystretch}{1.3}
\begin{tabularx}{\textwidth}{lXXXX}
\toprule
\textbf{Criterion} & \textbf{Elo} & \textbf{IRT 3PL} & \textbf{BKT} & \textbf{IRT+BKT Hybrid} \\
\midrule
Data requirements & Low & High & Medium & Medium (with priors) \\
Per-KC mastery tracking & No & No & Yes & Yes \\
Prerequisite enforcement & No & No & Yes & Yes \\
Item discrimination modelling & No & Yes & No & Yes \\
Knowledge decay detection & No & No & Partial & Yes (via decay flag) \\
Interpretability & High & High & High & High \\
Calibration stability (small $n$) & High & Low & Medium & Medium \\
\bottomrule
\end{tabularx}
\end{table}

% ─────────────────────────────────────────────────────────────────────────────
\section{Computerised Adaptive Testing}
\label{sec:lr_cat}
% ─────────────────────────────────────────────────────────────────────────────

CAT applies IRT to sequential item selection: each successive item is chosen
based on all previous responses rather than a fixed predetermined sequence
\parencite{wainer2000}. The foundational principle is maximum Fisher information
selection \parencite{weiss1982}:

\begin{equation}
  I(\theta) = \frac{a^2 \cdot P(\theta)(1-P(\theta))}{\left(\dfrac{P(\theta)-c}{1-c}\right)^2}
  \label{eq:fisher_information}
\end{equation}

Two methods exist for estimating $\theta$: Maximum Likelihood Estimation
(MLE) and Expected A Posteriori (EAP). MLE is asymptotically unbiased but
undefined when a learner answers all items correctly or incorrectly, which
is common in short sessions. EAP integrates the posterior over a discrete
grid of ability values:

\begin{equation}
  \hat{\theta}_{\text{EAP}} = \frac{\displaystyle\sum_k \theta_k \cdot L(\theta_k) \cdot \pi(\theta_k)}{\displaystyle\sum_k L(\theta_k) \cdot \pi(\theta_k)}
  \label{eq:eap}
\end{equation}

where $L(\theta_k)$ is the likelihood of the observed response pattern at grid
point $\theta_k$ and $\pi(\theta_k)$ is the prior probability. EAP always
produces a finite estimate and incorporates prior uncertainty, making it more
robust for short sessions \parencite{bock1982}. \textcite{baker2004} confirm
that EAP with a well-specified prior outperforms MLE on both bias and MSE
criteria in short-test conditions. This project implements EAP over 81 grid
points spanning $\theta \in [-4, 4]$ with a normal prior
$\theta_0 \sim \mathcal{N}(-0.780, 0.543)$.

% ─────────────────────────────────────────────────────────────────────────────
\section{Empirical Effectiveness}
\label{sec:lr_effectiveness}
% ─────────────────────────────────────────────────────────────────────────────

\textcite{vanlehn2011} synthesises 60 controlled studies, reporting a
hierarchy of tutoring modalities (Table~\ref{tab:effect_sizes}). The
near-equivalence of step-based ITS and human tutoring ($d = 0.76$ vs.\
$d = 0.79$) is the empirical justification for this system. \textcite{kulik2016}
report $d = 0.66$ for systems implementing continuous real-time adaptation
versus $d = 0.28$ for checkpoint systems, supporting the design decision to
implement per-response EAP updating. Geography presents an intermediate
epistemological structure: factual knowledge admits objective assessment;
procedural knowledge requires geographic frameworks; conceptual understanding
involves higher-order reasoning. \textcite{bednarz2013} identify significant
underrepresentation of UK Geography in adaptive learning research; the
predominance of North American and East Asian studies \parencite{matos2025}
means that the UK KS3/GCSE curriculum context remains insufficiently studied.

\begin{table}[H]
\centering
\caption{Effect size comparison across tutoring modalities \parencite{vanlehn2011}.}
\label{tab:effect_sizes}
\renewcommand{\arraystretch}{1.3}
\begin{tabular}{lc}
\toprule
\textbf{Tutoring modality} & \textbf{Effect size ($d$)} \\
\midrule
Human one-to-one tutoring & 0.79 \\
Step-based ITS & 0.76 \\
Answer-based tutoring systems & 0.31 \\
No-tutoring control & 0.00 (baseline) \\
\bottomrule
\end{tabular}
\end{table}

% ─────────────────────────────────────────────────────────────────────────────
\section{Implementation Challenges}
\label{sec:lr_challenges}
% ─────────────────────────────────────────────────────────────────────────────

\textbf{Cold-start problem.} Before sufficient learner data accumulates,
personalisation is necessarily suboptimal \parencite{matos2025}. This project
addresses cold-start through pre-test seeding: a 13-item diagnostic assessment
converts raw scores to a Laplace-smoothed logit estimate:

\begin{equation}
  \theta_{\text{seed}} = \frac{\ln\!\left(\dfrac{n_{\text{correct}} + 0.5}{n_{\text{incorrect}} + 0.5}\right)}{1.7}
  \label{eq:theta_seed}
\end{equation}

where 1.7 is the standard IRT logistic approximation scaling factor
\parencite{lord1980}. Laplace smoothing prevents boundary estimates at
perfect or zero scores.

\textbf{IRT calibration with small samples.} Traditional 3PL calibration
requires 500+ observations per item \parencite{lord1980}. The strategy adopted
is constrained estimation with $c = 0.25$ fixed, verified against pilot data
using RMSEA diagnostics, with a pre-specified Rasch fallback \parencite{embretson2000}.

\textbf{Usability.} \textcite{bangor2008} established SUS as the standard
educational technology usability metric; scores $\geq 85$ indicate excellent
usability, and adaptive systems typically score 15--20 points lower than
conventional interfaces. \textcite{mayer2009}'s multimedia design principles
(coherence, signalling, segmentation) directly inform the interface design:
the five-tab session summary provides coherence; the first-visit tour provides
signalling; pause-and-resume implements segmentation.

% ─────────────────────────────────────────────────────────────────────────────
\section{Ethical Considerations}
\label{sec:lr_ethics}
% ─────────────────────────────────────────────────────────────────────────────

\textcite{baker2021} identify three ethical dimensions for adaptive systems.
\textit{Transparency}: learners have a right to understand why the system
presents particular content; this system addresses this through the five-tab
session summary, making the adaptive engine's internal state visible and
interpretable. \textit{Data protection}: UK GDPR principles of data
minimisation, purpose limitation, and storage limitation \parencite{voigt2017}
are implemented through server-side consent flows, pseudonymised storage, and
JWT-based authentication compliant with \textcite{nist2017}.
\textit{Algorithmic bias}: adaptive systems can inadvertently reinforce
educational inequities through differential error rates across demographic
groups \parencite{baker2021}; this limitation is acknowledged as a constraint
on generalisability, with the convenience sample noted as non-representative
of the KS3/GCSE population.

% ─────────────────────────────────────────────────────────────────────────────
\section{Research Gaps and Project Contribution}
\label{sec:lr_gaps}
% ─────────────────────────────────────────────────────────────────────────────

The literature reviewed above reveals five significant gaps addressed by this
project.

\textbf{Gap 1: ITS application to UK Geography.}
No published ITS system targets UK KS3/GCSE Geography specifically. The
overwhelming majority of adaptive learning research addresses mathematics,
programming, and North American science curricula. This project provides
a proof-of-concept adaptive system for the UK Geography curriculum, extending
the empirical evidence base for ITS effectiveness beyond traditional STEM
domains.

\textbf{Gap 2: Hybrid IRT+BKT with open-source accessibility.}
While hybrid approaches are theoretically advocated \parencite{desmarais2012,
matos2025}, few open-source implementations combine IRT's item-level precision
with BKT's per-KC mastery tracking in a production-deployed web application.
This project provides a publicly accessible implementation at
\url{https://github.com/dlalloyd/E-Learning-Final-Project}.

\textbf{Gap 3: Bloom-gated adaptive progression.}
Existing ITS systems typically sequence items by IRT difficulty alone, without
enforcing cognitive-level prerequisites. This project implements Bloom-gated
selection as a live adaptive constraint, ensuring learners consolidate recall
before encountering analysis-level items.

\textbf{Gap 4: Long-term retention evaluation.}
The near-universal reliance on immediate post-test outcomes is a documented
methodological limitation \parencite{cepeda2006, vanlehn2011}. This project's
7-day delayed retention assessment provides evidence on the persistence of
adaptive learning gains.

\textbf{Gap 5: Interpretability under real-time adaptation.}
The interpretability--efficacy trade-off \parencite{matos2025}, in which more
powerful algorithms tend to be less interpretable, is addressed through the
combination of an interpretable probabilistic model (IRT+BKT) and a transparent
visualisation layer (the five-tab session summary).

% ─────────────────────────────────────────────────────────────────────────────
\section{Conclusions and Design Principles}
\label{sec:lr_conclusion}
% ─────────────────────────────────────────────────────────────────────────────

The literature establishes that adaptive ITS can approach human tutoring
effectiveness when designed with principled algorithm selection and transparent
feedback. The hybrid IRT+BKT approach is theoretically optimal for structured
multi-KC curricula with prerequisite dependencies. The following design
principles, derived from this review, inform the system architecture:

\begin{enumerate}[leftmargin=2em]
  \item \textbf{ZPD operationalisation}: target $P(\theta) \approx 0.70$ for item selection.
  \item \textbf{Mastery gating}: require $p(\text{Learned}) \geq 0.80$ per KC before progression.
  \item \textbf{Bloom-cognitive scaffolding}: enforce Bloom-level prerequisites as selection gates.
  \item \textbf{EAP estimation}: use Expected A Posteriori over a discretised grid for short-session stability.
  \item \textbf{Credit-adjusted BKT}: reduce the evidential weight of hint-assisted correct responses.
  \item \textbf{Successive relearning}: prioritise decayed KCs through discrimination-boosted selection.
  \item \textbf{Transparent adaptation}: surface internal model states through visualisation, not concealment.
\end{enumerate}
